\begin{figure}[!t]
\centering
\setlength{\unitlength}{1mm}
\begin{picture}(88,66)
\put(10,8){\framebox(70,50){}}
\put(10,8){\line(7,5){70}}

\put(10,8){\line(1,0){70}}
\put(10,8){\line(0,1){50}}

\put(10,8){\line(0,-1){1.5}}
\put(27.5,8){\line(0,-1){1.5}}
\put(45,8){\line(0,-1){1.5}}
\put(62.5,8){\line(0,-1){1.5}}
\put(80,8){\line(0,-1){1.5}}

\put(8.5,8){\line(-1,0){1.5}}
\put(8.5,20.5){\line(-1,0){1.5}}
\put(8.5,33){\line(-1,0){1.5}}
\put(8.5,45.5){\line(-1,0){1.5}}
\put(8.5,58){\line(-1,0){1.5}}

\put(9,3){\makebox(0,0){\scriptsize 0\%}}
\put(27.5,3){\makebox(0,0){\scriptsize 10\%}}
\put(45,3){\makebox(0,0){\scriptsize 20\%}}
\put(62.5,3){\makebox(0,0){\scriptsize 30\%}}
\put(80,3){\makebox(0,0){\scriptsize 40\%}}

\put(5,8){\makebox(0,0)[r]{\scriptsize 0\%}}
\put(5,20.5){\makebox(0,0)[r]{\scriptsize 10\%}}
\put(5,33){\makebox(0,0)[r]{\scriptsize 20\%}}
\put(5,45.5){\makebox(0,0)[r]{\scriptsize 30\%}}
\put(5,58){\makebox(0,0)[r]{\scriptsize 40\%}}

\put(45,0.5){\makebox(0,0){\footnotesize Population share}}
\put(1.5,33){\makebox(0,0){\rotatebox{90}{\footnotesize Normalized Banzhaf index}}}
\put(61,60.5){\makebox(0,0)[l]{\scriptsize equal-share line}}

\put(76.47,49.67){\circle*{1.8}}
\put(74.7,52.2){\makebox(0,0)[r]{\scriptsize Germany}}

\put(64.51,31.81){\circle*{1.8}}
\put(62.8,34.3){\makebox(0,0)[r]{\scriptsize France}}

\put(48.73,25.86){\circle*{1.8}}
\put(50.9,28.2){\makebox(0,0)[l]{\scriptsize Spain}}

\put(24.30,25.86){\circle*{1.8}}
\put(26.5,28.2){\makebox(0,0)[l]{\scriptsize Netherlands}}

\put(10.54,19.90){\circle*{1.8}}
\put(12.7,22.2){\makebox(0,0)[l]{\scriptsize Luxembourg}}

\put(10.45,19.90){\circle*{1.8}}
\put(12.7,17.6){\makebox(0,0)[l]{\scriptsize Malta}}
\end{picture}
\caption{Population share versus normalized Banzhaf power in the 6-country EU dual-majority mini-model. The figure shows that voting power does not move one-for-one with population share: Germany remains the most powerful player, but the Netherlands has substantially more Banzhaf power than its population share alone would suggest, while the largest states lie below the equal-share benchmark. This reflects the dual-majority structure, in which large states help satisfy the population threshold but each state contributes equally toward the state-count threshold.}
\label{fig:eu_mini_population_vs_banzhaf}
\end{figure}
